% ============================================================================
% CAPÍTULO 5: PARTÍCULAS ELEMENTALES
% ============================================================================
\chapter{Partículas Elementales}
\label{cap:particulas}

\begin{center}
\textit{``Los ladrillos del universo''}
\end{center}

\vspace{0.5cm}

\begin{tcolorbox}[colback=white, colframe=kleinblue, boxrule=2pt]
\centering
\textit{``¿Por qué el protón es 1836 veces más pesado que el electrón?''}\\
--- Pregunta sin respuesta, siglo XX\\[0.3cm]
\textit{``Porque $1836 = 6\pi^5$.''}\\
--- Teoría Klein, 2026
\end{tcolorbox}


% ============================================================================
\section{El Zoológico de Partículas}
% ============================================================================

\begin{analogiabox}
\textbf{El LEGO Cósmico}

Imagina que todo el universo está hecho de piezas de LEGO.
Pero no hay infinitos tipos de piezas --- solo unas pocas fundamentales.

\textbf{FERMIONES (materia):}

\begin{center}
\begin{tabular}{|c|c|}
\hline
\textbf{QUARKS} & \textbf{LEPTONES} \\
\hline
u (up), d (down) & e (electrón), $\nu_e$ (neutrino e) \\
c (charm), s (strange) & $\mu$ (muón), $\nu_\mu$ (neutrino $\mu$) \\
t (top), b (bottom) & $\tau$ (tau), $\nu_\tau$ (neutrino $\tau$) \\
\hline
\end{tabular}
\end{center}

\textbf{BOSONES (fuerzas):}
\begin{itemize}
    \item Fotón ($\gamma$): fuerza electromagnética
    \item Gluones (g): fuerza nuclear fuerte
    \item $W^\pm$, $Z^0$: fuerza nuclear débil
    \item Higgs (H): da masa a las partículas
    \item Gravitón (?): fuerza gravitacional (aún no detectado)
\end{itemize}
\end{analogiabox}

\textbf{EL MISTERIO DE LAS MASAS:}

¿Por qué estas partículas tienen las masas que tienen?
El Modelo Estándar NO lo explica --- son ``parámetros libres''.

\begin{center}
\begin{tabular}{ll}
Electrón: & $m_e = 0.511$ MeV \\
Muón: & $m_\mu = 105.66$ MeV \\
Tau: & $m_\tau = 1776.86$ MeV \\
Protón: & $m_p = 938.27$ MeV \\
Higgs: & $m_H = 125,250$ MeV \\
\end{tabular}
\end{center}

¿Hay un patrón? La Teoría Klein dice que SÍ.


% ============================================================================
\section{La Masa del Protón: $m_p/m_e = 6\pi^5$}
% ============================================================================

El ratio entre la masa del protón y el electrón es uno de los
números más importantes de la física.

\begin{nivelconceptual}
\textbf{¿Por qué importa $m_p/m_e$?}

Este ratio determina:
\begin{itemize}
    \item El tamaño de los átomos
    \item La velocidad de las reacciones químicas
    \item La estabilidad del hidrógeno
    \item Si la vida es posible
\end{itemize}

Si fuera muy diferente, no habría química, ni planetas, ni vida.

Valor observado: $m_p/m_e = 1836.15267$
\end{nivelconceptual}

\begin{nivelmatematico}
\textbf{La Fórmula Klein para $m_p/m_e$}

\begin{equation}
\boxed{\frac{m_p}{m_e} = 6\pi^5 = (7-1) \times \pi^5}
\end{equation}

\textbf{Desglose:}
\begin{itemize}
    \item $\pi^5 = 306.0197$ (quinta potencia de $\pi$)
    \item $6 = 7 - 1$ (capas Klein menos una de referencia)
    \item $6\pi^5 = 1836.1181$
\end{itemize}

\textbf{VERIFICACIÓN:}
\begin{align}
\text{Predicción:} & \quad 1836.1181\\
\text{Observado:} & \quad 1836.1527\\
\text{Error:} & \quad 0.0019\%
\end{align}

\textbf{¡ESTA ES NUESTRA MEJOR PREDICCIÓN!}\\
Solo 0.002\% de error --- mejor que muchas mediciones experimentales.
\end{nivelmatematico}

\begin{analogiabox}
\textbf{Las 6 Capas Activas}

La botella de Klein tiene 7 capas.
Pero UNA capa es especial: es la ``referencia'' o el ``origen''.

Las otras 6 capas son ``activas'' y contribuyen a la física.

\begin{verbatim}
     Capa 7 --+
     Capa 6   |
     Capa 5   | <-- 6 capas activas
     Capa 4   |
     Capa 3   |
     Capa 2 --+
     Capa 1 ---- <-- 1 capa de referencia
\end{verbatim}

El factor $6 = 7-1$ aparece en el ratio de masas porque
el protón ``siente'' las 6 capas activas, mientras el electrón
es una partícula fundamental que vive en la capa de referencia.
\end{analogiabox}

\textbf{¿POR QUÉ $\pi^5$?}

El exponente 5 no es arbitrario:
\begin{itemize}
    \item $5$ = dimensiones de Kaluza-Klein (4 espacio-tiempo + 1 extra)
    \item $5$ = número de dimensiones donde la botella de Klein es natural
    \item $5$ aparece en $\dim(\text{SU}(5)) = 24 = 5^2 - 1$
\end{itemize}

La masa del protón codifica la estructura 5-dimensional del universo.


% ============================================================================
\section{La Masa del Muón: $m_\mu/m_e = 21\pi^2 = 3 \times 7 \times \pi^2$}
% ============================================================================

El muón es como un ``electrón gordo'' --- misma carga, pero 207 veces más pesado.

\begin{nivelmatematico}
\textbf{La Fórmula Klein para $m_\mu/m_e$}

\begin{equation}
\boxed{\frac{m_\mu}{m_e} = 21\pi^2 = 3 \times 7 \times \pi^2}
\end{equation}

\textbf{Desglose:}
\begin{itemize}
    \item $\pi^2 = 9.8696$
    \item $21 = 3 \times 7$ (generaciones $\times$ capas Klein)
    \item $21\pi^2 = 207.26$
\end{itemize}

\textbf{VERIFICACIÓN:}
\begin{align}
\text{Predicción:} & \quad 207.26\\
\text{Observado:} & \quad 206.77\\
\text{Error:} & \quad 0.24\%
\end{align}
\end{nivelmatematico}

\textbf{INTERPRETACIÓN:}

El factor $21 = 3 \times 7$ tiene significado profundo:
\begin{itemize}
    \item $3$ = número de generaciones de fermiones (e, $\mu$, $\tau$)
    \item $7$ = número de capas de Klein
\end{itemize}

El muón es la ``segunda generación'' del electrón.
Su masa codifica tanto la estructura generacional (3) como la topológica (7).

Nota: $21 \approx 7\pi \approx 22$. ¡El muón también ``conoce'' la constante Klein!


% ============================================================================
\section{La Masa del Higgs: $m_H/m_p = 42.5\pi$}
% ============================================================================

El bosón de Higgs es especial: da masa a todas las demás partículas.

\begin{nivelconceptual}
\textbf{El Bosón de Higgs}

El campo de Higgs permea todo el universo como un ``jarabe cósmico''.
Las partículas que interactúan con él adquieren masa.

\begin{itemize}
    \item Fotones: no interactúan $\to$ sin masa
    \item Electrones: interactúan poco $\to$ masa pequeña
    \item Quarks top: interactúan mucho $\to$ masa enorme
\end{itemize}

El Higgs mismo tiene masa: $m_H = 125.25$ GeV
\end{nivelconceptual}

\begin{nivelmatematico}
\textbf{La Fórmula Klein para $m_H/m_p$}

\begin{equation}
\boxed{\frac{m_H}{m_p} = \left(6 \times 7 + \frac{1}{2}\right) \times \pi = 42.5\pi}
\end{equation}

\textbf{Desglose:}
\begin{itemize}
    \item $6 \times 7 = 42$ (capas activas $\times$ capas totales)
    \item $+ 1/2$ = corrección cuántica
    \item $\times \pi$ = factor geométrico
    \item $42.5\pi = 133.52$
\end{itemize}

\textbf{VERIFICACIÓN:}
\begin{align}
\text{Predicción:} & \quad 133.52\\
\text{Observado:} & \quad 133.50\\
\text{Error:} & \quad 0.02\%
\end{align}

¡Segunda mejor predicción después de $m_p/m_e$!
\end{nivelmatematico}

\begin{analogiabox}
\textbf{El Número 42}

En ``La Guía del Autoestopista Galáctico'', 42 es la respuesta
a la vida, el universo y todo lo demás.

¡Y resulta que $42 = 6 \times 7$ aparece en la masa del Higgs!

Douglas Adams quizás estaba más cerca de la verdad de lo que pensaba.

El Higgs, que da masa a todo, tiene un ratio de masa con el protón
de aproximadamente $42\pi$. El número 42 ESTÁ en la física fundamental.
\end{analogiabox}


% ============================================================================
\section{El Patrón de las Masas}
% ============================================================================

Reuniendo todas las fórmulas de masas:

\begin{table}[H]
\centering
\caption{Masas de Partículas desde Klein}
\begin{tabular}{lcccc}
\toprule
\textbf{Partícula} & \textbf{Ratio} & \textbf{Fórmula Klein} & \textbf{Error} & \textbf{Exp. de $\pi$} \\
\midrule
Protón & $m_p/m_e$ & $6\pi^5$ & 0.002\% & 5 \\
Higgs & $m_H/m_p$ & $42.5\pi$ & 0.02\% & 1 \\
Muón & $m_\mu/m_e$ & $21\pi^2$ & 0.24\% & 2 \\
\bottomrule
\end{tabular}
\end{table}

\textbf{PATRÓN EN LOS EXPONENTES: 1, 2, 5}

¡Son los primeros números de la secuencia de Fibonacci! (1, 1, 2, 3, 5, 8, 13, ...)

También son los primeros números primos: 2, 3, 5 (exceptuando el 1)

\begin{nivelconceptual}
\textbf{La Jerarquía de Masas}

Las masas de partículas varían enormemente:

\begin{center}
neutrino $<$ electrón $<$ muón $<$ protón $<$ Higgs $<$ top quark\\
$\sim$0.05 eV \quad 511 keV \quad 106 MeV \quad 938 MeV \quad 125 GeV \quad 173 GeV
\end{center}

Ratio extremo: $m_{\text{top}} / m_\nu \sim 10^{12}$

¿Por qué esta jerarquía? Klein responde:

\begin{itemize}
    \item Cada tipo de partícula ``siente'' diferente número de capas
    \item Las supresiones $(7\pi)^{-n}$ generan la jerarquía
    \item Los exponentes $n$ siguen patrones matemáticos (Fibonacci, primos)
\end{itemize}
\end{nivelconceptual}

\textbf{LOS COEFICIENTES (6, 21, 42.5):}

\begin{align}
6 &= 7 - 1 \quad \text{(capas activas)}\\
21 &= 3 \times 7 \quad \text{(generaciones} \times \text{capas)}\\
42.5 &= 6 \times 7 + 0.5 \quad \text{(activas} \times \text{totales + corrección)}
\end{align}

Todos involucran el número 7 (capas Klein).

El patrón es claro: las masas de partículas codifican la estructura
de 7 capas de la botella de Klein.


% ============================================================================
\section{Resumen del Capítulo}
% ============================================================================

\begin{tcolorbox}[colback=kleingray, colframe=kleinblue, title=\textbf{Resumen del Capítulo 5}]

\textbf{IDEAS CENTRALES:}

\begin{enumerate}
    \item Las masas de partículas elementales tienen fórmulas Klein simples.

    \item $m_p/m_e = 6\pi^5$ con 0.002\% de error (¡nuestra mejor predicción!)

    \item $m_H/m_p = 42.5\pi$ con 0.02\% de error (el número 42 es real)

    \item $m_\mu/m_e = 21\pi^2$ con 0.24\% de error

    \item Los exponentes de $\pi$ son 1, 2, 5 (Fibonacci/primos).
    Los coeficientes son 6, 21, 42.5 (todos involucran el 7).
\end{enumerate}

\vspace{0.3cm}
\textbf{ECUACIONES DEL CAPÍTULO:}

\begin{align}
m_p/m_e &= 6\pi^5 = (7-1)\pi^5\\
m_\mu/m_e &= 21\pi^2 = 3 \times 7 \times \pi^2\\
m_H/m_p &= 42.5\pi = (6 \times 7 + 1/2)\pi
\end{align}

\vspace{0.3cm}
\textbf{PRÓXIMO CAPÍTULO:}

¿Cómo explica Klein la cosmología: el CMB, la constante cosmológica,
y el número de Avogadro?

\end{tcolorbox}


% ============================================================================
\section*{Ejercicios}
% ============================================================================

\begin{enumerate}
    \item Verifica que $6\pi^5 \approx 1836$. ¿Cuál es el error exacto?

    \item El tau ($\tau$) es el leptón más pesado, con $m_\tau/m_e \approx 3477$.
    ¿Puedes encontrar una fórmula Klein para este número?
    (Pista: prueba con múltiplos de 7 y potencias de $\pi$)

    \item Si el protón fuera más liviano ($m_p/m_e = 1000$), ¿cómo cambiaría
    la química? ¿Serían los átomos más grandes o más pequeños?

    \item El número 42 aparece en la masa del Higgs. Investiga otros lugares
    en la física donde aparece 42 (o $6 \times 7$).
\end{enumerate}


% ============================================================================
\section*{Código de Verificación}
% ============================================================================

\begin{lstlisting}[caption={Verificación numérica - Capítulo 5}]
import numpy as np
from numpy import pi

# Proton/electron
mp_me_obs = 1836.15267
mp_me_klein = 6 * pi**5
print(f"PROTON/ELECTRON:")
print(f"  6*pi^5 = 6 * {pi**5:.4f} = {mp_me_klein:.5f}")
print(f"  Observado = {mp_me_obs:.5f}")
print(f"  Error = {abs(mp_me_klein - mp_me_obs)/mp_me_obs*100:.4f}%")

# Muon/electron
mu_me_obs = 206.7682830
mu_me_klein = 21 * pi**2
print(f"\nMUON/ELECTRON:")
print(f"  21*pi^2 = 21 * {pi**2:.6f} = {mu_me_klein:.4f}")
print(f"  Observado = {mu_me_obs:.4f}")
print(f"  Error = {abs(mu_me_klein - mu_me_obs)/mu_me_obs*100:.3f}%")

# Higgs/proton
mH_mp_obs = 125250 / 938.27
mH_mp_klein = 42.5 * pi
print(f"\nHIGGS/PROTON:")
print(f"  42.5*pi = 42.5 * {pi:.6f} = {mH_mp_klein:.4f}")
print(f"  Observado = {mH_mp_obs:.4f}")
print(f"  Error = {abs(mH_mp_klein - mH_mp_obs)/mH_mp_obs*100:.3f}%")

# Coeficientes y su relacion con 7
print(f"\nCOEFICIENTES Y SU RELACION CON 7:")
print(f"  6    = 7 - 1")
print(f"  21   = 3 * 7")
print(f"  42.5 = 6 * 7 + 0.5")
\end{lstlisting}
